Aggiungiamo al nucleo il meccanismo dei {\em checkpoint}.

Il checkpoint di un processo \`e la foto del suo stato in un certo istante, dato dal
contenuto dei registri del processore e contenuto di tutta la sua memoria privata.
Nel nostro caso la memoria privata comprende la pila sistema e la pila utente, ma
per evitare di dover rivelare il contenuto della pila sistema, permetteremo i checkpoint
solo negli istanti in cui un processo sta per ritornare a livello utente (e dunque la
sua pila sistema \`e vuota).

Nel meccanismo che vogliamo realizzare, un processo (master) si prepara a ricevere
i checkpoint da un altro processo (slave) tramite la primitiva \verb|cp_prep(pid)|.
La primitiva restituisce al master l'indirizzo a cui verr\`a salvato lo stato della
memoria privata (pila utente, per quanto detto sopra) dello slave ad ogni checkpoint.
Da quel momento in poi, ogni volta che lo slave sta per tornare a livello utente si
dovr\`a sincronizzare con il master per trasferigli un nuovo checkpoint. Il master
si pu\`o mettere in attesa del prossimo checkpoint tramite la primitiva \verb|cp_get(regs)|,
dove \verb|regs| \`e un puntatore ad un vettore di 16 \verb|natq|, destinato a ricevere
il contenuto dei registri dello slave. Dopo il trasferimento del checkpoint entrambi i
processi ritornano pronti. La relazione di master/slave prosegue fino a quando uno dei
due processi non termina. Ogni processo pu\`o essere al pi\`u master o slave di un altro processo.

Per realizzare il meccanismo appena descritto introduciamo la seguente struttura dati:
\begin{verbatim}
struct des_cp {
    natl master;
    natl slave;
    natq *cp_regs_master;
    bool slave_waiting;
    bool master_waiting;
};
\end{verbatim}
La struttura \verb|des_cp| descrive lo stato di una coppia master/slave.
Il campo \verb|master| \`e l'id del processo master;
il campo \verb|slave| \`e l'id del processo slave;
i campi \verb|slave_waiting| e \verb|master_waiting| valgono true se il corrispondente processo
\`e in attesa di trasferire/ricevere un nuovo checkpoint;
se il processo master \`e in attesa, il campo \verb|cp_regs_master| contiene il campo \verb|regs|
che il master ha passato a \verb|cp_get()|.

Aggiungiamo anche il seguente campo ai descrittori di processo:
\begin{verbatim}
    des_cp *cp;
\end{verbatim}
Nei processi che si trovano in relazione master/slave, questo campo punta ad un'istanza
di \verb|des_cp| (condivisa da entrambi i processi). Per i processi che non sono n\'e master,
n\'e slave, questo campo \`e nullo.

Aggiungiamo infine le seguenti primitive:
\begin{itemize}
  \item \verb|void* cp_prep(natl pid)| (gi\`a realizzata): registra il processo chiamante come master del processo di
    identificatore \verb|pid|, che diventa slave. Il processo chiamante non deve essere gi\`a master e lo
    slave deve essere un processo di livello utente diverso dal chiamante; restituisce \verb|nullptr| in caso
    di fallimento (il processo chiamante \`e slave, il processo \verb|pid| non esiste o \`e gi\`a slave o master,
    qualche allocazione di memoria \`e fallita), altrimenti restituisce l'indirizzo di una zona di memoria
    (nel master) pronta a contenere lo stato della memoria privata dello slave.
  \item \verb|bool cp_get(natq* regs)|: Il processo chiamante deve essere master. Il master si pone in attesa
    del prossimo checkpoint da parte dello slave. Quando la primitiva ritorna, \verb|regs| contiene lo stato
    dei registri, mentre lo stato della memoria privata si trova all'indirizzo precedentemente ricevuto
    tramite \verb|cp_prep()|. La primitiva restituisce \verb|false| se il processo slave \`e terminato
    prima di poter inviare un checkpoint.
\end{itemize}

Le primitive abortiscono il processo chiamante in caso di errore e tengono conto della priorit\`a tra i processi.
Si pu\`o assumere che la pila sistema occupi una singola pagina.

Modificare il file \verb|sistema.cpp| in modo da realizzare le parti mancanti.

{\bf Suggerimento}: per riconoscere quando un processo sta per tornare a livello utente ci si
pu\`o ispirare alla funzione \verb|liv_chiamante()|.
